\section{Conclusion}

We presented a systematic evaluation of PEFT methods for adapting Prithvi-100M to heterogeneous geospatial inputs. Across EuroSAT, BigEarthNet-S2, LandCover.ai, Linhe, and LoveDA, the results support a bounded conclusion: PEFT rankings familiar from NLP and conventional vision should be revalidated on remote sensing foundation models with architecture-aware controls. In the tested Prithvi-100M setting, Houlsby adapters provide the strongest and most reliable adaptation performance, and their advantage widens on dense prediction ($+155\%$ relative mIoU on LandCover.ai) rather than shrinking. Modality selection also has a large practical impact, especially when moving from reduced-band inputs to richer multi-spectral configurations.

The study also contributes a deeper diagnosis of LoRA failure in GeoFMs. Rather than a single negative result, we show a two-stage explanation: standard module-wise implementations can miss fused QKV projections, and corrected low-rank adaptation remains inferior to adapter-based methods in the tested Prithvi-100M settings. This finding highlights the importance of architecture-aware PEFT evaluation in remote sensing.

Our work has several limitations. We focus on a single backbone, Prithvi-100M, and use Esri-derived labels rather than independent manual labels for the Linhe production-style validation. Although the cross-dataset and cross-task consistency is strong, validation on additional geospatial foundation models, parameter-matched adapter sweeps, and independent Linhe label-quality checks are needed before the conclusions can be generalized beyond this backbone and deployment setting. Future work should also examine whether the same PEFT ranking holds for regression and temporal forecasting tasks.

Despite these limitations, the benchmark provides a useful empirical foundation for architecture-aware PEFT evaluation in remote sensing. By clarifying where standard LoRA insertion fails, where corrected low-rank adaptation remains insufficient, and where bottleneck adapters recover useful capacity, it offers both practical deployment guidance for Prithvi-100M and a methodological reference for future GeoFM adaptation studies.
